% Lab report in the Laboratory house style. Every element below shows the
% house convention for its kind: replace the pendulum with the experiment and
% keep the conventions. Build with build.py, which reads latexmkrc.
\documentclass[11pt]{article}
\usepackage{laboratory}

% The kind of document or the module code, set above the title; delete to omit.
\kicker{Laboratory report}
\title{Measuring the gravitational acceleration with a simple pendulum}
\author{Author Name}
\date{\today}
% The repository, notebook or data folder behind the report; delete to omit.
\source{analysis/pendulum.py}

\begin{document}

\maketitle

% Four sentences: aim, method, main result with its uncertainty, and what it means.
\begin{abstract}
We timed a simple pendulum at nine lengths between \qtylist{0.2;1.0}{\metre} and fitted the small-angle period to the measurements. The fit gives $g = \qty{9.809(4)}{\metre\per\second\squared}$, within one standard error of the local value of \qty{9.81}{\metre\per\second\squared}. The residuals show no trend with length, so the small-angle model describes the pendulum to the timing precision of \qty{1}{\milli\second}.
\end{abstract}

\section{Introduction}

% What is measured, why, and the model the measurement tests.
A mass on a light string of length $L$ swings with a period that depends only on $L$ and the gravitational acceleration $g$ when the amplitude is small~\cite{nelson1986}. Timing the period at several lengths therefore measures $g$, and the spread of the measurements about the model tests whether the small-angle approximation holds.

\section{Theory}

The angle $\theta$ of the string from the vertical obeys
\begin{equation}
  \frac{\dd^2\theta}{\dd t^2} = -\frac{g}{L} \sin\theta .
  \label{eq:motion}
\end{equation}
For small $\theta$, $\sin\theta \approx \theta$ and \cref{eq:motion} describes simple harmonic motion with period
\begin{equation}
  T = 2\pi \sqrt{L/g} .
  \label{eq:period}
\end{equation}
At an amplitude $\theta_0$ the true period is longer by a factor of about $1 + \theta_0^2/16$~\cite{nelson1986}, which is \num{1.7e-4} at the \qty{0.052}{\radian} used here, or \qty{0.3}{\milli\second} at the longest length, below the timing precision.

\section{Method}

% Apparatus, procedure and how each uncertainty was estimated.
A steel bob hung from a clamped string, released from rest at \qty{3}{\degree}. A light gate timed 20 oscillations at each length, and each period is the mean of five such runs with a standard error of \qty{1}{\milli\second}. Lengths were measured from the clamp to the centre of the bob with a metre rule to \qty{1}{\milli\metre}.

\section{Results}

% Lead with the result, then the evidence for it; every float is referred to before it appears.
\Cref{fig:period} shows the measured periods with the fit of \cref{eq:period}, and \cref{tab:period} lists them at selected lengths. A least-squares fit of $T = a \sqrt{L}$ gives $a = \qty{2.0061(4)}{\second\metre\tothe{-1/2}}$ and hence $g = 4\pi^2 / a^2 = \qty{9.809(4)}{\metre\per\second\squared}$. The fit has $\chi^2 = 6.9$ for 8 degrees of freedom.

\begin{figure}[tbp]
  \centering
  \includegraphics[width=\linewidth]{figures/pendulum}
  % The bold lead says what the figure shows; the keys name its marks in their colours.
  \caption{\lead{The small-angle model fits every length.} Period against length~\marker{seriesA}, with the fitted small-angle model~\linekey{seriesC} (a), and residuals from the fit (b). Error bars are one standard error.}
  \label{fig:period}
\end{figure}

% A table hangs from the left edge under its caption; it is not centred.
\begin{table}[tbp]
  \caption{Measured and fitted periods at selected lengths.}
  \label{tab:period}
  \begin{tabular}{
    S[table-format=1.1]
    S[table-format=1.4(2)]
    S[table-format=1.4]
    S[table-format=-1.1]
  }
    \toprule
    \multicolumn{2}{c}{\spanhead{Measured}} & \multicolumn{2}{c}{\spanhead{Fitted}} \\
    \cmidrule(lr){1-2} \cmidrule(lr){3-4}
    {\colhead{$L$}{\metre}} & {\colhead{$T$}{\second}} & {\colhead{$T_\mathrm{fit}$}{\second}} & {\colhead{$T - T_\mathrm{fit}$}{\milli\second}} \\
    \midrule
    0.2 & 0.8983(10) & 0.8972 & 1.1 \\
    0.4 & 1.2692(10) & 1.2688 & 0.4 \\
    0.6 & 1.5548(10) & 1.5540 & 0.8 \\
    0.8 & 1.7937(10) & 1.7943 & -0.6 \\
    1.0 & 2.0057(10) & 2.0061 & -0.4 \\
    \bottomrule
  \end{tabular}
\end{table}

\section{Discussion}

% What the result means, how far to trust it, and what would test it further.
The residuals scatter about zero with no trend in length, and a $\chi^2$ close to the number of degrees of freedom shows that the quoted timing error accounts for the scatter~\cite{taylor1997}. The fitted $g$ agrees with the local value, so neither the finite amplitude nor the mass of the string shifts the period by more than the timing precision. A constant error in length, such as a misjudged centre of the bob, changes the shape of the curve only slightly over this range, so the residuals test it weakly; measuring to the top and the bottom of the bob would bound it.

\section{Conclusion}

The simple pendulum gives $g = \qty{9.809(4)}{\metre\per\second\squared}$, and the small-angle model describes its period at every length measured.

\bibliographystyle{apsrev4-2}
\bibliography{refs}

\end{document}
